% STATUS: written 2026-11-05 (Kadison 1990 and Crist 1996 read in full)
% STATUS: extended 2026-11-06 (Ayupov, Kudaybergenov, Nurjanov and Alauatdinov 2011; Hadwin, Li, Li and Ma 2013; Arzikulov, Khakimov and Qurbonov 2026; Hao and Chen 2026 read in full)
% Sources read in full: Kadison 1990, J. Algebra 130, pp. 494--509; Crist 1996,
% J. Funct. Anal. 135, pp. 76--92; Ayupov, Kudaybergenov, Nurjanov and Alauatdinov 2011 (pp. 1--2, 7);
% Hadwin, Li, Li and Ma 2013 (pp. 1--3); Arzikulov, Khakimov and Qurbonov 2026 (pp. 1, 10--11, 15);
% Hao and Chen 2026 (pp. 1, 2, 4, 7).
% Excerpts in excerpts/Kadison1990_local.txt and excerpts/Crist1996_local.txt.
% Larson--Sourour 1990 could not be downloaded and is recorded as atlas-pending P0006.
\section{Local derivations: Kadison, Larson--Sourour}\label{sec:local-characterizations:local}

Chapter~\ref{ch:generalized-twisted} changed the derivation itself: it
twisted it, generalized it or gave it more arguments.
This chapter changes the form of the map
\emph{data}. Kadison names a linear map $\alpha$ of an algebra $R$ into an
$R$-bimodule $M$ a \emph{local derivation} in this case~\cite[p.~494]{Kadison1990Local}. At each $A\in R$
the value $\alpha(A)$ agrees with a true derivation $\delta_A$, whose choice
may vary with $A$ \cite[p.~494]{Kadison1990Local}. His paper opens the
study of such maps \cite[p.~494]{Kadison1990Local}.
Bre\v{s}ar and \v{S}emrl give the same definition in their abstract, and they
record that Kadison had already studied the question on von Neumann algebras
\cite[p.~483]{BresarSemrl1993}.

\begin{definition}[{\cite[p.~494]{Kadison1990Local}}]\label{def:local:local}
Let $R$ be an associative algebra and $M$ an $R$-bimodule. A linear mapping
$\alpha\colon R\to M$ is a \emph{local derivation} when for each $A\in R$
there is a derivation $\delta_{A}\colon R\to M$ such that
$\alpha(A)=\delta_{A}(A)$.
\end{definition}

The question is whether the varying witness can be replaced by one
derivation of Definition~\ref{def:local:local}. the criterion of Kadison answers
it for von Neumann algebras \cite[p.~494]{Kadison1990Local}. An $R$-bimodule $M$ is \emph{dual} when it
is a dual space and the module actions are weak* continuous
\cite[p.~494]{Kadison1990Local}.

\begin{proposition}[{\cite[Theorem~A, pp.~495--496]{Kadison1990Local}}]\label{prop:local:thmA}
Let $R$ be a von Neumann algebra and $M$ a dual $R$-bimodule. If
$\delta\colon R\to M$ is a norm-continuous local derivation, then $\delta$
is a derivation.
\end{proposition}

The proof is organized as five lemmas, and Lemma 4 asks for eleven
sublemmas \cite[p.~496]{Kadison1990Local}. The paragraph after the theorem
records the same reading of Proposition~\ref{prop:local:thmA}: under these
conditions each local derivation is a derivation \cite[p.~496]{Kadison1990Local}.

\begin{proposition}[{\cite[Theorem~B, p.~507]{Kadison1990Local}}]\label{prop:local:thmB}
Let $R$ be a von Neumann algebra. If $\delta\colon R\to R$ is a
norm-continuous linear mapping such that for each $A\in R$ there is a
$T_{A}\in R$ with $\delta(A)=AT_{A}-T_{A}A$, then there is a $T\in R$
with $\delta(A)=AT-TA$ for all $A\in R$.
\end{proposition}

Kadison reads Theorem B off Proposition~\ref{prop:local:thmA}. Every
derivation of a von Neumann algebra into itself is inner, so the map is a
local inner derivation and one inner derivation then serves all points
\cite[p.~507]{Kadison1990Local}. The statement may surprise even when $R$
is the full complex matrix algebra \cite[p.~494]{Kadison1990Local}.

A Banach-space companion came in the same year, 1990
\cite[p.~78]{Crist1996Local}. Larson and Sourour proved the same result for
$B(X)$, the bounded operators on a Banach space $X$
\cite[p.~78]{Crist1996Local}.
Bre\v{s}ar then lowered the hypothesis of Proposition~\ref{prop:local:thmA}.
A local derivation from a von Neumann algebra into a normed bimodule is
still a derivation \cite[p.~78]{Crist1996Local}.

Crist enlarges the stage from von Neumann algebras to operator algebras
with a maximal abelian subalgebra \cite[p.~76]{Crist1996Local}. A CSL
algebra is an operator algebra whose invariant subspaces form a commutative
lattice \cite[p.~76]{Crist1996Local}.

\begin{theorem}[{\cite[Theorem~2.7, p.~83]{Crist1996Local}}]\label{thm:local:csl}
Let $\mathcal{A}$ be a finite dimensional CSL algebra. Then any local
derivation $\varphi\colon\mathcal{A}\to\mathcal{A}$ must be a derivation.
\end{theorem}

The next theorem gives the same conclusion for inner local derivations.

\begin{theorem}[{\cite[Theorem~2.8, p.~88]{Crist1996Local}}]\label{thm:local:cslinner}
Let $\mathcal{A}$ be a finite dimensional CSL algebra. If
$\varphi\colon\mathcal{A}\to\mathcal{A}$ is an inner local derivation, then
$\varphi$ is an inner derivation.
\end{theorem}

In the language of the reflexivity program of Larson this says the derivations,
and relatively the inner derivations, are reflexive spaces on these
algebras \cite[pp.~78--79]{Crist1996Local}. The same conclusion passes to
algebras that are direct limits of finite dimensional CSL algebras by
$*$-extendable embeddings \cite[p.~78]{Crist1996Local}. One step is left
behind by Theorem~\ref{thm:local:csl}. Finite dimensional CSL algebras can
carry outer derivations, so the
results of Kadison and of Larson and Sourour do not carry over verbatim
\cite[p.~79]{Crist1996Local}.

The question of Larson for a general algebra stays open, and this section
states it. % claim-ok: own signpost introducing the moved problem, which carries its own citation

\begin{problem}[{\cite[p.~703]{HadwinLi2004}}]\label{prob:open:larson}
For which algebras $A$ is the space $\operatorname{der}(A)$ of bounded
derivations of $A$ reflexive, that is, equal to its own reflexive closure in the
bounded operators on $A$?
\end{problem}

Hadwin and Li answer the question for completely distributive commutative
subspace lattices and for $J$-subspace lattices, where the space of bounded
derivations is reflexive \cite[Theorem~2.15, p.~708]{HadwinLi2004}. They also
show that $\operatorname{der}(A,M)$ is algebraically $3$-reflexive for every
Banach $A$-bimodule $M$ \cite[Theorem~2.19, p.~710]{HadwinLi2004}. The local
property for automorphisms fails even in finite dimension, since
$T\mapsto AT^{t}A^{-1}$ on $M_n(\mathbb{C})$ is a surjective local automorphism
that is not an automorphism \cite[p.~713]{HadwinLi2004}.
The local condition also holds on algebras of measurable operators \cite[p.~1]{AyupovKudaybergenovNurjanovAlauatdinov2011}. Let $M$ be
a von Neumann algebra with a faithful normal semi-finite trace $\tau$, and let
$L^{\omega}(M,\tau)$ be the Arens algebra of all $\tau$-measurable operators
affiliated with $M$ \cite[p.~1]{AyupovKudaybergenovNurjanovAlauatdinov2011}.
Every local derivation of that algebra is a spatial derivation when $M$ is
finite \cite[Theorem~2.4, p.~7]{AyupovKudaybergenovNurjanovAlauatdinov2011}.
When $M$ is commutative the only local derivation of the Arens algebra is zero
\cite[Corollary~2.5, p.~7]{AyupovKudaybergenovNurjanovAlauatdinov2011}.

Hadwin, Li, Li and Ma answer a question of Kadison with a ring-theoretic
argument \cite[p.~1]{HadwinLiLiMa2013}. They prove the following. Let the abelian summand
of a von Neumann algebra $M$ be discrete. Then every local derivation of the
algebra of measurable operators affiliated with $M$ is a derivation
\cite[Theorem~1, pp.~2--3]{HadwinLiLiMa2013}. The proof passes through a lemma
on rings that carry a system of matrix units
\cite[Lemma~1, p.~1]{HadwinLiLiMa2013}.

On the two five-dimensional naturally graded nilpotent associative algebras
$\pi_2$ and $\pi_3$ the answer is negative \cite[p.~1]{ArzikulovKhakimovQurbonov2026}.
The matrix form of a derivation and the matrix form of a local derivation are
different on each algebra, so a local derivation need not be a derivation
\cite[p.~1]{ArzikulovKhakimovQurbonov2026}. The authors write both forms out on
$\pi_3$ \cite[Theorems~8 and~9, pp.~14--15]{ArzikulovKhakimovQurbonov2026}. The local
derivations of each algebra form a vector space, and that space is closed under
the commutator \cite[Theorems~10 and~11, p.~15]{ArzikulovKhakimovQurbonov2026}.
The matrix form of a local derivation obeys the rules that the matrix form of a
derivation obeys, and the authors list them
\cite[Theorem~5, pp.~10--11]{ArzikulovKhakimovQurbonov2026}.

The vector space of local derivations of an algebra is a Lie algebra under the
commutator, with no hypothesis on the algebra
\cite[Theorem~2.8, p.~4]{HaoChen2026Local}. Hao and Chen prove it for an
arbitrary algebra over a field, and their argument covers every space of
operators of one derivation type \cite[p.~4]{HaoChen2026Local}. The result
settles half of a question of Ayupov, Elduque and Kudaybergenov
\cite[p.~1]{HaoChen2026Local}. The other half concerns the local
automorphisms, and it is settled in the same paper
\cite[p.~1]{HaoChen2026Local}. The authors also prove that an algebra admits
an invertible local derivation exactly when it admits an invertible derivation
\cite[Theorem~2.18, p.~7]{HaoChen2026Local}.

\hbfigure{08-local-characterizations/local-pointwise}{A local derivation \cite[p.~494]{Kadison1990Local}. At each element $a$ of the algebra $R$ the value $\alpha(a)$ equals $\delta_{a}(a)$ for some derivation $\delta_{a}$ of $R$, and the derivation may change with $a$. The question the section answers is when one derivation serves every point.}{fig:local:pointwise}

Liu shows that every 2-local Lie derivation of a semi-finite factor von Neumann algebra is a Lie derivation~\cite{WlW2LocalLieDerivationsOnSemiFini2016}.

Niazi and Miri prove that every continuous local triple derivation from a $C^*$-algebra into any of its iterated duals, taken with the ternary module structure the authors introduced, is a triple derivation, a partial answer to a problem of Peralta and Russo~\cite{WlLocalTripleDerivationsFromCSta2018}.
